\begin{code}[hide]
module slides.Tsinghua-PL-seminar.sections.design2 where

open import Data.Nat using (ℕ)
open import Data.Product using (_×_ ; _,_)
open import Data.Sum using (_⊎_)
open import Data.List using (List)
open import Data.Unit using (⊤)
open import Level using (0ℓ)
import Relation.Binary.PropositionalEquality as Eq

open import Word.Base using (Word ; WRel ; [_]ʷ ; ε ; _•_)
open import Presentation.Construct.Base using ([_]ₗ ; [_]ᵣ)

infix 5 _⋄_⋄_

postulate
  A B : Set
  u v : Word A
\end{code}
%% Part 4, continued — a real coset level, products, stdlib (paper §4).
\begin{frame}{A real coset level, in full: adjoining $X$ to $\langle\omega\rangle\times\langle S\rangle$ \InPaper}
\begin{code}[hide]
module Sω where
  postulate ω S : Word ⊤
module CT where
  Sω = ⊤
  data XSω : Set where
    X-gen S-gen ω-gen : XSω
  data Cr : Set where
    X-cr ε-cr : Cr
  infixr 8 _^_
  postulate
    _^_ : ∀ {A : Set} → Word A → ℕ → Word A
    _===′_ : WRel XSω
    _~′_ : Word Sω × Cr → Word Sω × Cr → Set
    by-equal-nf′ : ∀ {A : Set} → A
  open Eq using (refl)
  data _==_ : WRel XSω where
    order-ω order-S order-X order-SX : ε == ε
    comm : ∀ {gen : XSω} → ε == ε
\end{code}
\begin{columns}[T]
\begin{column}{0.47\textwidth}
\small Base $\langle\omega\rangle\times\langle S\rangle$ (radices $8\cdot 4$);
adjoin $X$ with $X^2=\varepsilon$, $(SX)^2=\omega^2$, $\omega$ central.
Two cosets, $\varepsilon$ and $X$. The whole table:
\begin{code}
  h : Cr → XSω → Word Sω × Cr
  h X-cr X-gen = ε , ε-cr
  h X-cr S-gen = (Sω.ω ^ 2 • Sω.S ^ 3) , X-cr
  h X-cr ω-gen = Sω.ω , X-cr
  h ε-cr X-gen = ε , X-cr
  h ε-cr S-gen = Sω.S , ε-cr
  h ε-cr ω-gen = Sω.ω , ε-cr
\end{code}
{\scriptsize from \texttt{Examples/Amalgamations/CliffordT1.agda} (coset type renamed \aid{Cr})}
\end{column}
\begin{column}{0.50\textwidth}
\small Hypothesis (2), \emph{every} case closed by evaluation:
\begin{code}
  h-wd-ax : ∀ c {u t} → u == t → Set
  h-wd-ax X-cr order-ω = by-equal-nf′ refl , refl
  h-wd-ax X-cr order-S = by-equal-nf′ refl , refl
  h-wd-ax X-cr order-X = by-equal-nf′ refl , refl
  h-wd-ax X-cr order-SX = by-equal-nf′ refl , refl
  h-wd-ax X-cr (comm {X-gen}) = by-equal-nf′ refl , refl
  h-wd-ax X-cr (comm {S-gen}) = by-equal-nf′ refl , refl
  h-wd-ax X-cr (comm {ω-gen}) = by-equal-nf′ refl , refl
  -- … and the same seven for ε-cr
\end{code}
{\scriptsize abridged: \aid{(by-equal-nf Eq.refl) , Eq.refl} in the source}
\end{column}
\end{columns}
\begin{keybox}
\small The typechecker runs the table on both sides of the axiom: the cosets
coincide and the residues have the same normal form one level down.
Qubit Clifford+$T$: \textbf{59} such obligations; qutrit's nine-coset level
alone: \textbf{72} --- every one \aid{refl}.
\end{keybox}
\end{frame}

\begin{frame}{The one case that is a real derivation \InPaper}
\begin{columns}[c]
\begin{column}{0.62\textwidth}
\small Hypothesis (5) at coset $X$, letter $S$: the table says
$X\cdot S = (\omega^2 S^3)\cdot X$. The proof, verbatim:
\vspace{2pt}

{\agdafont\scriptsize
\begin{tabular}{@{}l@{}}
\textcolor{AgdaFunction}{lemma-XS} : \textcolor{AgdaFunction}{X} \textcolor{AgdaInductiveConstructor}{•} \textcolor{AgdaFunction}{S} \textcolor{AgdaDatatype}{≈} (\textcolor{AgdaFunction}{ω} \textcolor{AgdaFunction}{\textasciicircum} 2 \textcolor{AgdaInductiveConstructor}{•} \textcolor{AgdaFunction}{S} \textcolor{AgdaFunction}{\textasciicircum} 3) \textcolor{AgdaInductiveConstructor}{•} \textcolor{AgdaFunction}{X}\\
\textcolor{AgdaFunction}{lemma-XS} = \textcolor{AgdaFunction}{begin}\\
\quad \textcolor{AgdaFunction}{X} \textcolor{AgdaInductiveConstructor}{•} \textcolor{AgdaFunction}{S} \textcolor{AgdaFunction}{≈⟨} \textcolor{AgdaInductiveConstructor}{trans} (\textcolor{AgdaInductiveConstructor}{sym} \textcolor{AgdaInductiveConstructor}{right-unit}) (\textcolor{AgdaInductiveConstructor}{sym} (\textcolor{AgdaInductiveConstructor}{cong} \textcolor{AgdaInductiveConstructor}{refl} (\textcolor{AgdaInductiveConstructor}{axiom} \textcolor{AgdaInductiveConstructor}{order-X}))) \textcolor{AgdaFunction}{⟩}\\
\quad (\textcolor{AgdaFunction}{X} \textcolor{AgdaInductiveConstructor}{•} \textcolor{AgdaFunction}{S}) \textcolor{AgdaInductiveConstructor}{•} (\textcolor{AgdaFunction}{X} \textcolor{AgdaInductiveConstructor}{•} \textcolor{AgdaFunction}{X}) \textcolor{AgdaFunction}{≈⟨} \textcolor{AgdaInductiveConstructor}{trans} (\textcolor{AgdaInductiveConstructor}{sym} \textcolor{AgdaInductiveConstructor}{left-unit}) (\textcolor{AgdaInductiveConstructor}{sym} (\textcolor{AgdaInductiveConstructor}{cong} (\textcolor{AgdaInductiveConstructor}{axiom} \textcolor{AgdaInductiveConstructor}{order-S}) \textcolor{AgdaInductiveConstructor}{refl})) \textcolor{AgdaFunction}{⟩}\\
\quad (\textcolor{AgdaFunction}{S} \textcolor{AgdaFunction}{\textasciicircum} 4) \textcolor{AgdaInductiveConstructor}{•} (\textcolor{AgdaFunction}{X} \textcolor{AgdaInductiveConstructor}{•} \textcolor{AgdaFunction}{S}) \textcolor{AgdaInductiveConstructor}{•} (\textcolor{AgdaFunction}{X} \textcolor{AgdaInductiveConstructor}{•} \textcolor{AgdaFunction}{X}) \textcolor{AgdaFunction}{≈⟨} \textcolor{AgdaFunction}{by-assoc} \textcolor{AgdaInductiveConstructor}{Eq.refl} \textcolor{AgdaFunction}{⟩}\\
\quad (\textcolor{AgdaFunction}{S} \textcolor{AgdaFunction}{\textasciicircum} 3) \textcolor{AgdaInductiveConstructor}{•} (\textcolor{AgdaFunction}{S} \textcolor{AgdaInductiveConstructor}{•} \textcolor{AgdaFunction}{X}) \textcolor{AgdaFunction}{\textasciicircum} 2 \textcolor{AgdaInductiveConstructor}{•} \textcolor{AgdaFunction}{X} \textcolor{AgdaFunction}{≈⟨} \textcolor{AgdaInductiveConstructor}{cong} \textcolor{AgdaInductiveConstructor}{refl} (\textcolor{AgdaInductiveConstructor}{cong} (\textcolor{AgdaInductiveConstructor}{axiom} \textcolor{AgdaInductiveConstructor}{order-SX}) \textcolor{AgdaInductiveConstructor}{refl}) \textcolor{AgdaFunction}{⟩}\\
\quad (\textcolor{AgdaFunction}{S} \textcolor{AgdaFunction}{\textasciicircum} 3) \textcolor{AgdaInductiveConstructor}{•} \textcolor{AgdaFunction}{ω} \textcolor{AgdaFunction}{\textasciicircum} 2 \textcolor{AgdaInductiveConstructor}{•} \textcolor{AgdaFunction}{X} \textcolor{AgdaFunction}{≈⟨} \textcolor{AgdaInductiveConstructor}{sym} \textcolor{AgdaInductiveConstructor}{assoc} \textcolor{AgdaFunction}{⟩}\\
\quad (\textcolor{AgdaFunction}{S} \textcolor{AgdaFunction}{\textasciicircum} 3 \textcolor{AgdaInductiveConstructor}{•} \textcolor{AgdaFunction}{ω} \textcolor{AgdaFunction}{\textasciicircum} 2) \textcolor{AgdaInductiveConstructor}{•} \textcolor{AgdaFunction}{X} \textcolor{AgdaFunction}{≈⟨} \textcolor{AgdaInductiveConstructor}{cong} (\textcolor{AgdaFunction}{lemma-ω\textasciicircum{}n} 2 (\textcolor{AgdaFunction}{S} \textcolor{AgdaFunction}{\textasciicircum} 3)) \textcolor{AgdaInductiveConstructor}{refl} \textcolor{AgdaFunction}{⟩}\\
\quad (\textcolor{AgdaFunction}{ω} \textcolor{AgdaFunction}{\textasciicircum} 2 \textcolor{AgdaInductiveConstructor}{•} \textcolor{AgdaFunction}{S} \textcolor{AgdaFunction}{\textasciicircum} 3) \textcolor{AgdaInductiveConstructor}{•} \textcolor{AgdaFunction}{X} \textcolor{AgdaFunction}{∎}
\end{tabular}}
\end{column}
\begin{column}{0.35\textwidth}
\begin{keybox}[title=Division of labour]
\small
\begin{itemize}
  \item \textbf{coverage checker}: no case is missed;
  \item \textbf{evaluator}: closes the cases that are computation;
  \item \textbf{human}: writes the few real derivations, which read
        line by line like the pen-and-paper ones.
\end{itemize}
\end{keybox}
\vspace{4pt}
{\scriptsize (typeset by hand in Agda's colours: this block needs the full
module context of \texttt{CliffordT1.agda})}
\end{column}
\end{columns}
\end{frame}

\begin{frame}{Products of presentations \InPaper}
\begin{columns}[c]
\begin{column}{0.56\textwidth}
Relation sets over $A \uplus B$ from one primitive: a three-way join whose
\alert{mixed} component is the design parameter.
\begin{code}
data _⋄_⋄_ (Γ₁ : WRel A) (Γ₂ : WRel B) (Γ₃ : WRel (A ⊎ B))
           : WRel (A ⊎ B) where
  left  : Γ₁ u v → (Γ₁ ⋄ Γ₂ ⋄ Γ₃) [ u ]ₗ [ v ]ₗ
  right : Γ₂ u v → (Γ₁ ⋄ Γ₂ ⋄ Γ₃) [ u ]ᵣ [ v ]ᵣ
  mid   : Γ₃ u v → (Γ₁ ⋄ Γ₂ ⋄ Γ₃) u v
\end{code}
\vspace{2pt}
{\small Congruence-lifting lemmas (\aid{lefts}, \aid{rights}) embed the factors;
normal forms transport across.}
\end{column}
\begin{column}{0.41\textwidth}
\small
\renewcommand{\arraystretch}{1.25}
\begin{tabular}{@{}ll@{}}
\toprule
mixed relation $\Gamma_3$ & construction\\
\midrule
\aid{EmptyRel} & free product\\
\aid{CommRel} & direct product $\times$\\
\aid{ConjRel}, \aid{ConjRelʷ} & semidirect $\rtimes$\\
\aid{AmalgRel} $f_1\,f_2$ & amalgamated $\ast_M$\\
\aid{SugarRel} & definitional extension\\
$\Gamma$ \aid{⊕\textasciicircum} $n$ & $n$-fold power\\
\bottomrule
\end{tabular}

\vspace{6pt}
{\scriptsize \aid{ConjRelʷ}: moving $h$ past $n$ conjugates $n$ by a
\emph{word} --- \aid{conj : H → N → Word N}.}
\end{column}
\end{columns}
\end{frame}

\begin{frame}{Constructions come in pairs \InPaper}
\begin{columns}[c]
\begin{column}{0.50\textwidth}
\textbf{Syntactic}: normal-form lifting.
\begin{itemize}
  \item direct / semidirect: carrier $\mathit{NF}_1\times\mathit{NF}_2$;
  \item amalgam: a normal form for the \emph{common subgroup} plus one verified
        coset table per factor:
\end{itemize}
\begin{code}[hide]
amalNFC : Set → Set → Set
\end{code}
\begin{code}
amalNFC C D = (D ⊎ ⊤) × List (C × D) × (C ⊎ ⊤)
\end{code}
\textbf{Semantic}: presentation theorems against concrete groups
\[
\begin{array}{l}
(\Gamma \aidop{⋄} \Delta \aidop{⋄} \aid{CommRel})\ \aid{IsPresentationOf}\ (G_1 \times G_2) \\[2pt]
(\Gamma \aidop{⋄} \Delta \aidop{⋄} \aid{ConjRelʷ}\;\mathit{conj})\ \aid{IsPresentationOf}\ (G_1 \rtimes G_2) \\[2pt]
(\Gamma \mathrel{\aid{⊕\textasciicircum}} n)\ \aid{IsPresentationOf}\ (\otimes\aid{-group}\;n) \\[2pt]
(P_1 \mathrel{\aid{*}} P_2 \mathrel{\aid{⋆}} f_1 \mathrel{\aid{⋆}} f_2)\ \aid{IsPresentationOf}\ (G_1 \ast_{G_0} G_2)
\end{array}
\]
\end{column}
\begin{column}{0.47\textwidth}
\centering
{\small alternating normal form of an amalgam:}\\[4pt]
\begin{tikzpicture}[x=8.5mm,
  b0/.style={draw=gray,fill=gray!15,minimum height=7mm,minimum width=8mm,font=\scriptsize},
  bd/.style={draw=paperC,fill=softpaper,minimum height=7mm,minimum width=8mm,font=\scriptsize},
  bc/.style={draw=accent,fill=soft,minimum height=7mm,minimum width=8mm,font=\scriptsize}]
  \node[b0] at (0,0) {$\mathcal{M}$};
  \node[bd] at (1,0) {$d_0$};
  \node[bc] at (2,0) {$c_1$};
  \node[bd] at (3,0) {$d_1$};
  \node at (4,0) {$\cdots$};
  \node[bc] at (5,0) {$c_k$};
  \node[bd] at (6,0) {$d_k$};
  \node[bc] at (7,0) {$c$};
  \draw[decorate,decoration={brace,amplitude=3pt,mirror}] (0.55,-0.5) -- (1.45,-0.5) node[midway,below=2pt,font=\tiny]{$D\uplus\top$};
  \draw[decorate,decoration={brace,amplitude=3pt,mirror}] (1.55,-0.5) -- (6.45,-0.5) node[midway,below=2pt,font=\tiny]{\aid{List (C × D)}};
  \draw[decorate,decoration={brace,amplitude=3pt,mirror}] (6.55,-0.5) -- (7.45,-0.5) node[midway,below=2pt,font=\tiny]{$C\uplus\top$};
  \draw[decorate,decoration={brace,amplitude=3pt}] (-0.45,0.5) -- (0.45,0.5) node[midway,above=2pt,font=\tiny]{$\mathit{NF}_0$};
\end{tikzpicture}

\vspace{8pt}
{\small Group \alert{extensions}: given $1\to N\to G\to Q\to 1$ with the
conjugation action and the corrections by which lifted relators of $Q$ fail,
$N$'s relations + conjugation rules + corrected lifts present $G$.}

\vspace{4pt}
\begin{tikzpicture}[>={Stealth[length=1.6mm]},x=11mm]
  \foreach \x/\l in {0/1,1/N,2/G,3/Q,4/1} \node (n\x) at (\x,0) {$\l$};
  \foreach \a/\b in {0/1,1/2,2/3,3/4} \draw[->] (n\a) -- (n\b);
  \draw[->,dashed,newC] (n3) to[bend right=40] node[above,font=\scriptsize]{lift} (n2);
\end{tikzpicture}
\end{column}
\end{columns}
\end{frame}

\begin{frame}{The semantic side: contributions to the standard library \InPaper}
\begin{columns}[c]
\begin{column}{0.55\textwidth}
A presentation theorem cannot be stated without the group it presents --- and
stdlib had the structures but not the \alert{constructions}.
\begin{itemize}
  \setlength{\itemsep}{2pt}
  \item \texttt{Algebra.Construct.SemiDirectProduct}: $N\rtimes H$ via a
        bundled \aid{Action} (six laws);
  \item \texttt{Algebra.Construct.Amalgamation}: $M\ast_C N$ of
        \emph{monoids}, as a setoid --- the pushout;
  \item \texttt{Algebra.Construct.Extension}, \texttt{.SplitExtension}:
        short exact sequences;
  \item \texttt{Algebra.Morphism.Consequences}: monoid homomorphisms between
        groups are group homomorphisms;
  \item \aid{IsMonoidEpimorphism};\; \aid{Permutation′ n} as a \aid{Group}.
\end{itemize}
\small Seven modules, 865 lines --- being prepared for upstreaming.
\end{column}
\begin{column}{0.42\textwidth}
\begin{keybox}[title={Solvers by evaluation, not reflection}]
\small
\begin{itemize}
  \item \aid{by-assoc}: same letter sequence $\Rightarrow$ $w\approx v$;
  \item \aid{by-passoc}: re-bracket toward a pattern with holes;
  \item \aid{by-equal-nf}: same computed normal form.
\end{itemize}
Each maps both sides into a canonical data type by a function proved
correct once, then compares by \aid{refl}: the evaluator does the work of
a reflective tactic.
\end{keybox}
\end{column}
\end{columns}
\end{frame}
